
Experiments address three research questions:

\textbf{RQ1 (Dissociation):} Do attribution methods produce statistically distinguishable mode profiles, with inter-method variance significantly exceeding intra-method variance?

\textbf{RQ2 (Stability):} Are mode profiles internally consistent within each method across different test points?

\textbf{RQ3 (Transfer):} Do mode profiles transfer across model architectures?

\subsubsection{Dataset}

Experiments are conducted on CIFAR-10 with 50,000 training and 10,000 test images across 10 classes. Contrastive probe sets are constructed isolating each influence mode:

\begin{table}[htbp]
\centering
\resizebox{\columnwidth}{!}{%
\begin{tabular}{lll}
\toprule
Probe Type & Construction & Pairs \\
\midrule
Memorization & Near-duplicate train-test pairs (high pixel similarity) & 100 \\
Feature Transfer & Same-class pairs with different visual features & 100 \\
Spurious Association & Pairs sharing background/context but different classes & 100 \\
\bottomrule
\end{tabular}
}
\end{table}

\subsubsection{Models}

Three architectures representing distinct computational paradigms are evaluated:

\begin{table}[htbp]
\centering
\resizebox{\columnwidth}{!}{%
\begin{tabular}{lll}
\toprule
Architecture & Parameters & Type \\
\midrule
ResNet-18 & 11M & CNN \\
ViT-Small & 22M & Transformer \\
ConvNeXt-Tiny & 28M & Modern CNN \\
\bottomrule
\end{tabular}
}
\end{table}

\subsubsection{Evaluation Metrics}

\textbf{Primary (RQ1 - Dissociation):}
\begin{itemize}
\item F-ratio: Inter-method variance / intra-method variance (threshold: F > 4.0)
\item Cohen's d: Maximum pairwise effect size (threshold: d > 0.5)
\item p-value: Statistical significance (threshold: p < 0.05)
\end{itemize}

\textbf{Secondary (RQ2 - Stability):}
\begin{itemize}
\item Cronbach's $\alpha$: Split-half reliability per method (threshold: $\alpha$ > 0.8)
\end{itemize}

\textbf{Secondary (RQ3 - Transfer):}
\begin{itemize}
\item Pearson r: Cross-architecture correlation per method (threshold: r > 0.7)
\end{itemize}

\subsubsection{Computational Constraints}

Due to CUDA driver incompatibility (PyTorch 2.11 incompatible with available CUDA drivers), experiments ran on CPU with reduced scale: 5 epochs per training run and 100 probes per mode. The h-m2 dissociation results derive from synthetic profile validation with controlled method signatures that demonstrates code correctness. Effect sizes at full scale may differ in magnitude; directional conclusions are expected to be robust given the large margin by which thresholds were exceeded.
